%% Owner: Víctor Martín Castilla
\section{Conclusions}
\label{sec:conclusions}

The basic and app stacks run on a three-node Swarm on Azure, behind one
public entry with HTTPS, and the atenea stack is built, tested on a
laptop and ready to deploy. \todo{Mikael: one sentence on the stats
stack.} Data outlives containers: the app's database moved between nodes
with its data intact, and Moodle came back from a backup in
\qty{21}{\second}. Secrets never reach the public repository, stateless
services update one replica at a time with automatic rollback, and the
internet reaches only ports 22, 80 and 443 of one node.

Several things do not work as they should yet:
\begin{itemize}
  \item The registry login is merged but not active. On 19 October the
        registry still answered without a password, so anyone could push
        an image under a tag the cluster runs.
  \item The backups share a disk with the data they protect
        (\cref{sec:data}).
  \item nginx runs once, pinned to the node that keeps its certificates,
        and it does not reload them after a renewal.
  \item All members share one administrator account.
  \item The logs of most containers never rotate.
\end{itemize}

Three lessons stand out. First, the details of shared storage decide
what containers can do: root squashing, a sensible default, broke the
registry, and turning it off made every container that mounts the share
root over all of it. Second, Swarm has no start order, so a stack must
coordinate its own start, as Moodle's install marker does. Third, a
backup only counts after a timed restore: Moodle's first drill seemed to
fail, and only the event log showed why.

The obvious next step is an off-site copy of the backup folders, which
would keep the recovery objectives even if \texttt{swarm-manager-1} were
lost. After it come the registry login, the certificates on the shared
disk so nginx can scale, and personal accounts for each member.
